\uuid{BY9w}
\exo7id{3839}
\titre{exo7 3839}
\auteur{quercia}
\organisation{exo7}
\datecreate{2010-03-11}
\isIndication{false}
\isCorrection{false}
\chapitre{Espace euclidien, espace normé}
\sousChapitre{Espaces vectoriels hermitiens}
\cpge{mpsi-mp2i.espaces-vectoriels.a, pcsi.espaces-vectoriels.a, ptsi.geometrie-espace.c, mpsi-mp2i.espaces-vectoriels.b, pcsi.espaces-vectoriels.b, ptsi.espaces-vectoriels.a, mp-mpi.topologie-espaces-normes.a, pc.espaces-vectoriels-normes.a, psi.espaces-vectoriels-normes.a, mp-mpi.topologie-espaces-normes.b, pc.espaces-vectoriels-normes.b, psi.espaces-vectoriels-normes.b, psi.espaces-prehilbertiens.b, pt.structure-prehilbertienne.c}
\module{Algèbre}
\niveau{L2}
\difficulte{}

\contenu{
\texte{
Soit $E = \mathcal{C}([0,1],\R)$ muni du produit scalaire usuel, $F$ le sous-espace vectoriel des fonctions
polynomiales et $g$ la fonction exponentielle sur $[0,1]$.
}
\begin{enumerate}
    \item \question{Montrer que $g\notin F$.}
    \item \question{Montrer qu'il existe une suite $(f_n)$ de fonctions polynomiales convergeant
    vers $g$ pour la norme eucli\-dienne.}
    \item \question{En déduire que $F$ n'a pas de supplémentaire orthogonal.}
\end{enumerate}
}
